\uuid{TbRd}
\exo7id{1484}
\titre{exo7 1484}
\auteur{ridde}
\organisation{exo7}
\datecreate{1999-11-01}
\isIndication{false}
\isCorrection{false}
\chapitre{Espace euclidien, espace normé}
\sousChapitre{Projection, symétrie}
\cpge{mpsi-mp2i.matrices-applications-lineaires.a, pcsi.matrices-applications-lineaires.a, ptsi.matrices-applications-lineaires.a, mpsi-mp2i.matrices-applications-lineaires.b, pcsi.matrices-applications-lineaires.b, ptsi.matrices-applications-lineaires.b, mpsi-mp2i.changements-bases.c, mp-mpi.endomorphismes-euclidien.b, pc.endomorphismes-euclidiens.b, psi.endomorphismes-euclidien.b, pt.isometries-euclidien.b, mp-mpi.endomorphismes-euclidien.e, psi.endomorphismes-euclidien.e, pt.isometries-euclidien.d, bcpst2.applications-lineaires-matrices.c, bcpst2.valeurs-propres-vecteurs-propres.a, bcpst2.produit-scalaire.c}
\module{Algèbre}
\niveau{L2}
\difficulte{}

\contenu{
\question{
Quelle est la transformation de $\Rr^{3}$ dont la matrice dans la base canonique
est $\frac{1}{7} \begin{pmatrix} -2 & 6&-3\\ 6&3&2\\-3&2&6\\ \end{pmatrix}$ ?
}
}
